% ------------------------------------------------------------------
% Filter-diagonalization section (ready to include in a paper)
% ------------------------------------------------------------------
\section{Polynomial Filter Construction and Application}
\label{sec:filter}

We seek eigenpairs of a Hamiltonian operator $H$ in a target energy region.
The filtering stage is designed to (i) isolate selected energy windows and
(ii) do so with a fixed number of Hamiltonian applications independent of the
number of windows.

\subsection{Step 1: Energy interval selection and spectral scaling}

Let the physical spectral interval of interest be
\begin{equation}
E \in [E_{\min},E_{\max}], \qquad
V_{\min}=E_{\min}, \qquad dE = E_{\max}-E_{\min}.
\end{equation}
We define the scaled operator
\begin{equation}
\widetilde{H} = \frac{4}{dE}(H - V_{\min}I) - 2I,
\label{eq:Hscale}
\end{equation}
so that the spectrum is mapped to (approximately) $[-2,2]$.

This scaling is numerically important: after repeated polynomial recursions
(e.g., $\sim 500$ iterations), intermediate vectors remain in a safe dynamic
range for double precision because the effective argument of the polynomial is
confined near $[-2,2]$.

\subsection{Step 2: Gaussian filter width from interpolation order}

For a window centered at $E_\ell$, the Gaussian filter is
\begin{equation}
f_\ell(E)=\sqrt{\frac{dt}{\pi}}\exp\!\left[-dt\,(E-E_\ell)^2\right],
\end{equation}
with bandwidth
\begin{equation}
\sigma = \frac{1}{\sqrt{2dt}}.
\end{equation}
In our implementation, $dt$ is chosen from the interpolation order $n$ (named
\texttt{nc} in code) and global energy span $dE$ via
\begin{equation}
dt = \left(\frac{n}{2.5\,dE}\right)^2,
\qquad
\sigma = \frac{1}{\sqrt{2dt}}.
\label{eq:dt_sigma}
\end{equation}
Hence increasing $n$ yields narrower Gaussian windows (smaller $\sigma$), while
also improving interpolation accuracy.

\subsection{Step 3: Newton interpolation of each filter}

Choose interpolation nodes $\{s_j\}_{j=0}^{n-1}\subset[-2,2]$ in scaled
coordinates and map them back to physical energies by
\begin{equation}
E(s)=\frac{(s+2)dE}{4}+V_{\min}.
\end{equation}
For each window center $E_\ell$, compute Newton divided-difference
coefficients $a_{\ell,j}$ such that
\begin{equation}
P_\ell(s)=\sum_{j=0}^{n-1} a_{\ell,j}
\prod_{k=0}^{j-1}(s-s_k)
\approx f_\ell(E(s)).
\label{eq:newton_poly}
\end{equation}
Then $P_\ell(\widetilde{H})$ approximates $f_\ell(H)$.

\subsection{Step 4: Shared recursion for $N_{\mathrm{windows}}$ filters}

A direct strategy would apply each filter independently, costing
$\mathcal{O}(N_{\mathrm{windows}}\,n)$ Hamiltonian applications.
Instead, we share one Newton basis recursion among all windows.

Define basis vectors for an input state $\psi$:
\begin{align}
\phi_0 &= \psi,\\
\phi_j &= (\widetilde{H}-s_{j-1}I)\phi_{j-1}, \qquad j=1,\dots,n-1.
\label{eq:basis_rec}
\end{align}
Each filtered state is a linear combination
\begin{equation}
\psi^{(\ell)}_{\mathrm{filt}} = \sum_{j=0}^{n-1} a_{\ell,j}\,\phi_j.
\label{eq:combine}
\end{equation}
Therefore, all $N_{\mathrm{windows}}$ outputs are produced with only
$n-1$ applications of $H$ total (equivalently $n$ Newton levels), not
$N_{\mathrm{windows}}\times n$.

\subsection{Operator application with Horner/Newton recursion}

The implementation evaluates the Newton polynomial action via nested
recurrence (Horner form for Newton basis), while accumulating all windows in
parallel through different coefficient rows $\{a_{\ell,j}\}_{\ell}$.

\begin{algorithm}[t]
\caption{Shared Newton--Horner Filter Application for Multiple Windows}
\label{alg:shared_filter}
\begin{algorithmic}[1]
\Require Hamiltonian apply routine $\mathrm{ApplyH}(\cdot)$, initial state $\psi$,
         nodes $s[0{:}n-1]$, coefficients $a[0{:}N_w-1,0{:}n-1]$, parameters $(V_{\min},dE)$
\Ensure Filtered states $\Psi_{\mathrm{out}}[\ell] \approx f_\ell(H)\psi$, $\ell=0,\dots,N_w-1$
\State $\phi \gets \psi$
\State $\Psi_{\mathrm{out}}[\ell] \gets a[\ell,0]\,\psi$ for all $\ell$
\For{$j=1$ to $n-1$}
    \State $u \gets \mathrm{ApplyH}(\phi)$
    \State $\phi \gets \dfrac{4}{dE}(u - V_{\min}\phi) - 2\phi - s[j-1]\phi$
    \For{$\ell=0$ to $N_w-1$}
        \State $\Psi_{\mathrm{out}}[\ell] \gets \Psi_{\mathrm{out}}[\ell] + a[\ell,j]\,\phi$
    \EndFor
\EndFor
\State \Return $\Psi_{\mathrm{out}}$
\end{algorithmic}
\end{algorithm}

\appendix
\section{Node selection and node-ordering strategy for Newton interpolation}
\label{app:nodes}

This appendix summarizes the node-construction logic used in the code for
Newton interpolation. The key implementation detail is that both
\\emph{node locations} and \\emph{node order} are treated as first-class
numerical design choices.

\subsection{Why node order matters}

For nodes $\{s_j\}_{j=0}^{n-1}$, the Newton form is
\begin{equation}
P(s)=a_0+a_1(s-s_0)+a_2(s-s_0)(s-s_1)+\cdots .
\end{equation}
The divided-difference coefficients $\{a_j\}$ depend on the \\emph{ordered}
sequence $(s_0,s_1,\ldots,s_{n-1})$, not only on the unordered set.
Therefore, reordering nodes (e.g., sorting by coordinate) changes the
high-order coefficients and can significantly alter numerical stability.
In this implementation, when nodes are generated by greedy procedures
(Ashkenazy/Leja), their greedy order is preserved.

\subsection{Base node sets on the scaled interval $[-2,2]$}

All nodes are built in scaled coordinates $s\in[-2,2]$.
The code supports four base strategies:
\begin{enumerate}
  \item \textbf{Ashkenazy greedy nodes.}
  On a dense candidate pool, choose nodes sequentially by maximizing the
  incremental log-Vandermonde objective:
  \begin{equation}
  s_j=\arg\max_{s\in\mathcal C}
  \sum_{i=0}^{j-1}\log|s-s_i|^2.
  \label{eq:ashkenazy_obj}
  \end{equation}
  This produces an endpoint-enriched distribution similar to Chebyshev-like
  behavior while keeping a stable greedy order.

  \item \textbf{Chebyshev nodes.}
  First-kind Chebyshev nodes mapped to $[-2,2]$:
  \begin{equation}
  s_k = \frac{S_{\min}+S_{\max}}{2}
      + \frac{S_{\max}-S_{\min}}{2}
      \cos\!\left(\frac{(2k+1)\pi}{2n}\right),
  \end{equation}
  with $S_{\min}=-2$, $S_{\max}=2$.

  \item \textbf{Derivative-adapted nodes (for bandpass-like windows).}
  Build a sampling density from transition-region derivatives, then sample
  inverse-CDF quantiles to place more nodes near steep filter edges.

  \item \textbf{Density-mapped nodes.}
  Define a custom density $\rho(s)$ (typically Gaussian-enhanced in a target
  sub-interval), compute the normalized cumulative map, run Ashkenazy in the
  mapped coordinate, then invert the map back to $s$.
\end{enumerate}

\subsection{Adaptive refinement: Leja extension without moving old nodes}

When local interpolation accuracy in a specified energy interval is
insufficient, the code does \\emph{not} regenerate all nodes from scratch.
Instead, it appends new nodes with a Leja-style extension:
\begin{equation}
s_{n+j}
=\arg\max_{s\in\mathcal C}
\sum_{i=0}^{n+j-1}\log|s-s_i|^2,
\qquad j=0,1,\dots
\label{eq:leja_append}
\end{equation}
while freezing existing nodes $\{s_0,\dots,s_{n-1}\}$.
This preserves the previously fitted polynomial structure and avoids global
node drift.

Let $I_{\mathrm{enh}}=[s_{\mathrm{lo}},s_{\mathrm{hi}}]$ be the enhancement
interval (mapped from physical energy bounds), and let
$\varepsilon_{\max}$ be the maximum interpolation error in $I_{\mathrm{enh}}$
over all windows:
\begin{equation}
\varepsilon_{\max}
= \max_{\ell}\max_{s\in I_{\mathrm{enh}}}
\left|f_\ell(E(s)) - P_\ell(s)\right|.
\end{equation}
Refinement repeats in batches until
$\varepsilon_{\max}\le \tau_{\mathrm{interp}}$ or a maximum number of
iterations is reached.

\begin{algorithm}[t]
\caption{Node generation and ordering for stable Newton interpolation}
\label{alg:node_strategy}
\begin{algorithmic}[1]
\Require Initial order $n$, method in \{Ashkenazy, Chebyshev, derivative-adapted, density-mapped\}, optional enhancement interval
\Ensure Ordered node sequence $s[0{:}n_{\mathrm{true}}-1]$ and coefficients $a_{\ell,j}$
\State Build base nodes $s[0{:}n-1]$ in scaled domain $[-2,2]$ using the selected method
\State Keep generated order unchanged (no coordinate sorting)
\State Compute Newton coefficients $a_{\ell,:}$ for each window center $E_\ell$
\If{enhancement interval is provided}
  \State Evaluate worst interpolation error $\varepsilon_{\max}$ in the interval
  \While{$\varepsilon_{\max}>\tau_{\mathrm{interp}}$ and iteration limit not reached}
    \State Append nodes by Leja objective in Eq.~\eqref{eq:leja_append} (old nodes fixed)
    \State Recompute coefficients $a_{\ell,:}$ on the extended ordered node list
    \State Re-evaluate $\varepsilon_{\max}$
  \EndWhile
\EndIf
\State \Return final ordered nodes and Newton coefficients
\end{algorithmic}
\end{algorithm}

\paragraph{Complexity note.}
Algorithm~\ref{alg:shared_filter} requires only $n-1$ calls to
$\mathrm{ApplyH}$ for all windows together; additional windows only increase
cheap axpy-style vector combinations.

\subsection{Ritz extraction after filtering: SVD and projected diagonalization}

Let $\{\psi^{(i)}_{\mathrm{filt}}\}_{i=1}^{N_b}$ denote all filtered states
(from all random seeds and all windows), reshaped as columns of
\begin{equation}
B = [\psi^{(1)}_{\mathrm{filt}},\ldots,\psi^{(N_b)}_{\mathrm{filt}}]
\in \mathbb{C}^{N_g\times N_b},
\end{equation}
where $N_g$ is the full grid/basis dimension.

Because filtered vectors are often nearly linearly dependent, we first build a
numerically stable orthonormal subspace by QR+SVD:
\begin{align}
B &= QR,\\
R &= U_1 \Sigma V^\ast,\\
U_r &= (Q U_1)_{(:,1:r)},
\end{align}
where $r$ is the effective rank selected by a singular-value threshold
$\sigma_i>\tau_{\mathrm{svd}}$. The columns of $U_r\in\mathbb{C}^{N_g\times r}$
form the reduced Ritz subspace.

We then construct the projected (Rayleigh--Ritz) matrix
\begin{equation}
\widetilde{H}_{ij}
= \langle u_i, H u_j\rangle,
\qquad i,j=1,\ldots,r,
\label{eq:Htilde}
\end{equation}
or in matrix form
\begin{equation}
\widetilde{H}=U_r^\ast H U_r\in\mathbb{C}^{r\times r}.
\end{equation}
For Hermitian problems, we optionally enforce numerical symmetry by
\begin{equation}
\widetilde{H}\leftarrow \frac{\widetilde{H}+\widetilde{H}^\ast}{2}.
\end{equation}

Finally, solve the reduced eigenproblem
\begin{equation}
\widetilde{H} y_k = \theta_k y_k,
\qquad k=1,\ldots,r.
\end{equation}
The scalars $\theta_k$ are Ritz values (approximations to eigenvalues of $H$),
and the Ritz vectors in the full space are
\begin{equation}
x_k = U_r y_k.
\end{equation}
Sorting $\theta_k$ in ascending order gives the final energy estimates in the
target window; the corresponding $x_k$ are the approximate eigenstates.

\begin{algorithm}[t]
\caption{SVD-accelerated Rayleigh--Ritz from filtered basis}
\label{alg:ritz}
\begin{algorithmic}[1]
\Require Filtered basis matrix $B\in\mathbb{C}^{N_g\times N_b}$, operator apply $\mathrm{ApplyH}(\cdot)$, SVD cutoff $\tau_{\mathrm{svd}}$
\Ensure Ritz pairs $(\theta_k, x_k)$
\State Remove zero/invalid columns of $B$ and normalize remaining columns
\State Compute reduced QR: $B=QR$
\State SVD of $R$: $R=U_1\Sigma V^\ast$
\State $r \gets \#\{i:\Sigma_{ii}>\tau_{\mathrm{svd}}\}$, then $r\gets \max(1,r)$
\State $U_r \gets (Q U_1)_{(:,1:r)}$
\State Build $\widetilde{H}\in\mathbb{C}^{r\times r}$ by
      $\widetilde{H}_{ij}=\langle U_r[:,i], \mathrm{ApplyH}(U_r[:,j])\rangle$
\State (Optional Hermitian cleanup) $\widetilde{H}\gets (\widetilde{H}+\widetilde{H}^\ast)/2$
\State Diagonalize: $\widetilde{H}Y=Y\Theta$
\State Ritz values: $\theta_k=\Theta_{kk}$; Ritz vectors: $x_k=U_rY[:,k]$
\State Sort $(\theta_k,x_k)$ by ascending $\theta_k$
\State \Return $\{(\theta_k,x_k)\}$
\end{algorithmic}
\end{algorithm}
